%DIF 1-2c1-2
%DIF LATEXDIFF DIFFERENCE FILE
%DIF DEL initial_submission/main.tex                Tue Aug 18 10:58:26 2026
%DIF ADD revision/manuscript_clean/manuscript.tex   Tue Aug 18 19:09:19 2026
%DIF < % Template for PLoS
%DIF < % Version 3.8 Apr 2026
%DIF -------
% PLOS ONE manuscript -- PONE-D-26-28965 (revision) %DIF > 
% Based on the official PLOS LaTeX template v3.8 (Apr 2026). %DIF > 
%DIF -------
%
%DIF 4c4
%DIF < % % % % % % % % % % % % % % % % % % % % % %
%DIF -------
% % % % % % % % % % % % % % % % % % % % % % % %DIF > 
%DIF -------
%
%DIF 6c6
%DIF < % -- IMPORTANT NOTE
%DIF -------
% THIS IS THE CLEAN MANUSCRIPT. %DIF > 
%DIF -------
%
%DIF 8-11c8-11
%DIF < % This template contains comments intended
%DIF < % to minimize problems and delays during our production
%DIF < % process. Please follow the template instructions
%DIF < % whenever possible.
%DIF -------
% It contains no revision highlighting and no embedded graphics; PLOS requires %DIF > 
% the clean "Manuscript" file to be free of tracked changes and requires each %DIF > 
% figure to be uploaded as an individual file separate from the manuscript %DIF > 
% (https://journals.plos.org/plosone/s/figures). %DIF > 
%DIF -------
%
%DIF 13c13-15
%DIF < % % % % % % % % % % % % % % % % % % % % % % %
%DIF -------
% The "Revised Manuscript with Track Changes" file is produced from this source %DIF > 
% with latexdiff against the previously submitted version; see ../README.md. %DIF > 
% Do not hand-edit colors into this file. %DIF > 
%DIF -------
%
%DIF 15-20d17
%DIF < % Once your paper is accepted for publication,
%DIF < % PLEASE REMOVE ALL TRACKED CHANGES in this file
%DIF < % and leave only the final text of your manuscript.
%DIF < % PLOS recommends the use of latexdiff to track changes during review, as this will help to maintain a clean tex file.
%DIF < % Visit https://www.ctan.org/pkg/latexdiff?lang=en for info.
%DIF < %
%DIF -------
% % % % % % % % % % % % % % % % % % % % % % %
%
%DIF 23-29c19-24
%DIF < % NOTE ON REVISION HIGHLIGHTING
%DIF < % Passages added or modified in response to reviewer comments are highlighted
%DIF < % in yellow using the \rev{...} command (text) and \cellcolor{yellow} (table
%DIF < % cells). This highlighting is for the review stage only and should be removed
%DIF < % before final/production submission (PLOS asks that the final text contain no
%DIF < % in-text colors). To strip all highlighting at once, redefine the macro as
%DIF < % \renewcommand{\rev}[1]{#1} in the preamble.
%DIF -------
% NOTE ON \rev / \revs %DIF > 
% The wrappers \rev{...} and \revs{...} mark passages that were added or %DIF > 
% modified in the first revision round. They are deliberately defined as %DIF > 
% identity macros so that they never render color, while preserving the record %DIF > 
% of which passages changed. Keep them; do not re-enable the colored versions %DIF > 
% in this file. %DIF > 
%DIF -------
%
% % % % % % % % % % % % % % % % % % % % % % %

\documentclass[10pt,letterpaper]{article}
\usepackage[top=0.85in,left=2.75in,footskip=0.75in]{geometry}

% amsmath and amssymb packages, useful for mathematical formulas and symbols
\usepackage{amsmath,amssymb}

% Use adjustwidth environment to exceed column width (see example table in text)
\usepackage{changepage}

% textcomp package and marvosym package for additional characters
\usepackage{textcomp,marvosym}

% cite package, to clean up citations in the main text. Do not remove.
\usepackage{cite}

% Use nameref to cite supporting information files (see Supporting Information section for more info)
\usepackage{nameref,hyperref}

% line numbers
\usepackage[right]{lineno}

% ligatures disabled
\usepackage[nopatch=eqnum]{microtype}
\DisableLigatures[f]{encoding = *, family = * }

% color can be used to apply background shading to table cells only
\usepackage[table]{xcolor}

% array package and thick rules for tables
\usepackage{array}

%DIF 64-73c59-64
%DIF < % Revision highlighting (yellow) using xcolor only -- no extra packages needed.
%DIF < % \rev{...}  : block-level highlight for added/changed sentences or paragraphs
%DIF < % \revs{...} : short inline highlight (e.g. inside a heading), single line
%DIF < % To remove all highlighting before final submission, redefine both as identity:
%DIF < %   \renewcommand{\rev}[1]{#1}   \renewcommand{\revs}[1]{#1}
%DIF < \newcommand{\rev}[1]{\par\noindent\colorbox{yellow}{\parbox{\dimexpr\linewidth-2\fboxsep\relax}{#1}}\par}
%DIF < \newcommand{\revs}[1]{\colorbox{yellow}{#1}}
%DIF < % --- highlighting disabled for final/clean version ---
%DIF < \renewcommand{\rev}[1]{#1}
%DIF < \renewcommand{\revs}[1]{#1}
%DIF -------
% Round-1 revision markers, defined as identity macros so that the clean %DIF > 
% manuscript contains no in-text color (PLOS ONE requirement for the %DIF > 
% unmarked "Manuscript" file). \rev{...} wraps changed paragraphs and %DIF > 
% \revs{...} wraps changed inline fragments such as headings. %DIF > 
\newcommand{\rev}[1]{#1} %DIF > 
\newcommand{\revs}[1]{#1} %DIF > 
%DIF -------

% create "+" rule type for thick vertical lines
\newcolumntype{+}{!{\vrule width 2pt}}

% create \thickcline for thick horizontal lines of variable length
\newlength\savedwidth
\newcommand\thickcline[1]{%
  \noalign{\global\savedwidth\arrayrulewidth\global\arrayrulewidth 2pt}%
  \cline{#1}%
  \noalign{\vskip\arrayrulewidth}%
  \noalign{\global\arrayrulewidth\savedwidth}%
}

% \thickhline command for thick horizontal lines that span the table
\newcommand\thickhline{\noalign{\global\savedwidth\arrayrulewidth\global\arrayrulewidth 2pt}%
\hline
\noalign{\global\arrayrulewidth\savedwidth}}


%DIF 93-95c84-86
%DIF < % Remove comment for double spacing
%DIF < %\usepackage{setspace}
%DIF < %\doublespacing
%DIF -------
% PLOS ONE requires double-spaced manuscript text. %DIF > 
\usepackage{setspace} %DIF > 
\doublespacing %DIF > 
%DIF -------

% Text layout
\raggedright
\setlength{\parindent}{0.5cm}
\textwidth 5.25in
\textheight 8.75in

% Bold the 'Figure #' in the caption and separate it from the title/caption with a period
% Captions will be left justified
\usepackage[aboveskip=1pt,labelfont=bf,labelsep=period,justification=raggedright,singlelinecheck=off]{caption}
\renewcommand{\figurename}{Fig}

% Please use the included `plos2025.bst` as your BibTeX style
\bibliographystyle{plos2025}

% Remove brackets from numbering in List of References
\makeatletter
\renewcommand{\@biblabel}[1]{\quad#1.}
\makeatother



% Header and Footer with logo
\usepackage{lastpage,fancyhdr,graphicx}
\usepackage{epstopdf}
%\pagestyle{myheadings}
\pagestyle{fancy}
\fancyhf{}
%\setlength{\headheight}{27.023pt}
%\lhead{\includegraphics[width=2.0in]{PLOS-submission.eps}}
\rfoot{\thepage/\pageref{LastPage}}
\renewcommand{\headrulewidth}{0pt}
\renewcommand{\footrule}{\hrule height 2pt \vspace{2mm}}
\fancyheadoffset[L]{2.25in}
\fancyfootoffset[L]{2.25in}
\lfoot{\today}

%% Include all user-defined macros below

%% END MACROS SECTION
%DIF PREAMBLE EXTENSION ADDED BY LATEXDIFF
%DIF UNDERLINE PREAMBLE %DIF PREAMBLE
\RequirePackage[normalem]{ulem} %DIF PREAMBLE
\RequirePackage{color}\definecolor{RED}{rgb}{1,0,0}\definecolor{BLUE}{rgb}{0,0,1} %DIF PREAMBLE
\providecommand{\DIFaddtex}[1]{{\protect\color{blue}\uwave{#1}}} %DIF PREAMBLE
\providecommand{\DIFdeltex}[1]{{\protect\color{red}\sout{#1}}} %DIF PREAMBLE
%DIF SAFE PREAMBLE %DIF PREAMBLE
\providecommand{\DIFaddbegin}{} %DIF PREAMBLE
\providecommand{\DIFaddend}{} %DIF PREAMBLE
\providecommand{\DIFdelbegin}{} %DIF PREAMBLE
\providecommand{\DIFdelend}{} %DIF PREAMBLE
\providecommand{\DIFmodbegin}{} %DIF PREAMBLE
\providecommand{\DIFmodend}{} %DIF PREAMBLE
%DIF FLOATSAFE PREAMBLE %DIF PREAMBLE
\providecommand{\DIFaddFL}[1]{\DIFadd{#1}} %DIF PREAMBLE
\providecommand{\DIFdelFL}[1]{\DIFdel{#1}} %DIF PREAMBLE
\providecommand{\DIFaddbeginFL}{} %DIF PREAMBLE
\providecommand{\DIFaddendFL}{} %DIF PREAMBLE
\providecommand{\DIFdelbeginFL}{} %DIF PREAMBLE
\providecommand{\DIFdelendFL}{} %DIF PREAMBLE
%DIF HYPERREF PREAMBLE %DIF PREAMBLE
\providecommand{\DIFadd}[1]{\texorpdfstring{\DIFaddtex{#1}}{#1}} %DIF PREAMBLE
\providecommand{\DIFdel}[1]{\texorpdfstring{\DIFdeltex{#1}}{}} %DIF PREAMBLE
\providecommand{\DIFscaledelfig}{0.5}
%DIF HIGHLIGHTGRAPHICS PREAMBLE %DIF PREAMBLE
\RequirePackage{settobox} %DIF PREAMBLE
\RequirePackage{letltxmacro} %DIF PREAMBLE
\newsavebox{\DIFdelgraphicsbox} %DIF PREAMBLE
\newlength{\DIFdelgraphicswidth} %DIF PREAMBLE
\newlength{\DIFdelgraphicsheight} %DIF PREAMBLE
% store original definition of \includegraphics %DIF PREAMBLE
\LetLtxMacro{\DIFOincludegraphics}{\includegraphics} %DIF PREAMBLE
\providecommand{\DIFaddincludegraphics}[2][]{{\color{blue}\fbox{\DIFOincludegraphics[#1]{#2}}}} %DIF PREAMBLE
\providecommand{\DIFdelincludegraphics}[2][]{% %DIF PREAMBLE
\sbox{\DIFdelgraphicsbox}{\DIFOincludegraphics[#1]{#2}}% %DIF PREAMBLE
\settoboxwidth{\DIFdelgraphicswidth}{\DIFdelgraphicsbox} %DIF PREAMBLE
\settoboxtotalheight{\DIFdelgraphicsheight}{\DIFdelgraphicsbox} %DIF PREAMBLE
\scalebox{\DIFscaledelfig}{% %DIF PREAMBLE
\parbox[b]{\DIFdelgraphicswidth}{\usebox{\DIFdelgraphicsbox}\\[-\baselineskip] \rule{\DIFdelgraphicswidth}{0em}}\llap{\resizebox{\DIFdelgraphicswidth}{\DIFdelgraphicsheight}{% %DIF PREAMBLE
\setlength{\unitlength}{\DIFdelgraphicswidth}% %DIF PREAMBLE
\begin{picture}(1,1)% %DIF PREAMBLE
\thicklines\linethickness{2pt} %DIF PREAMBLE
{\color[rgb]{1,0,0}\put(0,0){\framebox(1,1){}}}% %DIF PREAMBLE
{\color[rgb]{1,0,0}\put(0,0){\line( 1,1){1}}}% %DIF PREAMBLE
{\color[rgb]{1,0,0}\put(0,1){\line(1,-1){1}}}% %DIF PREAMBLE
\end{picture}% %DIF PREAMBLE
}\hspace*{3pt}}} %DIF PREAMBLE
} %DIF PREAMBLE
\LetLtxMacro{\DIFOaddbegin}{\DIFaddbegin} %DIF PREAMBLE
\LetLtxMacro{\DIFOaddend}{\DIFaddend} %DIF PREAMBLE
\LetLtxMacro{\DIFOdelbegin}{\DIFdelbegin} %DIF PREAMBLE
\LetLtxMacro{\DIFOdelend}{\DIFdelend} %DIF PREAMBLE
\DeclareRobustCommand{\DIFaddbegin}{\DIFOaddbegin \let\includegraphics\DIFaddincludegraphics} %DIF PREAMBLE
\DeclareRobustCommand{\DIFaddend}{\DIFOaddend \let\includegraphics\DIFOincludegraphics} %DIF PREAMBLE
\DeclareRobustCommand{\DIFdelbegin}{\DIFOdelbegin \let\includegraphics\DIFdelincludegraphics} %DIF PREAMBLE
\DeclareRobustCommand{\DIFdelend}{\DIFOaddend \let\includegraphics\DIFOincludegraphics} %DIF PREAMBLE
\LetLtxMacro{\DIFOaddbeginFL}{\DIFaddbeginFL} %DIF PREAMBLE
\LetLtxMacro{\DIFOaddendFL}{\DIFaddendFL} %DIF PREAMBLE
\LetLtxMacro{\DIFOdelbeginFL}{\DIFdelbeginFL} %DIF PREAMBLE
\LetLtxMacro{\DIFOdelendFL}{\DIFdelendFL} %DIF PREAMBLE
\DeclareRobustCommand{\DIFaddbeginFL}{\DIFOaddbeginFL \let\includegraphics\DIFaddincludegraphics} %DIF PREAMBLE
\DeclareRobustCommand{\DIFaddendFL}{\DIFOaddendFL \let\includegraphics\DIFOincludegraphics} %DIF PREAMBLE
\DeclareRobustCommand{\DIFdelbeginFL}{\DIFOdelbeginFL \let\includegraphics\DIFdelincludegraphics} %DIF PREAMBLE
\DeclareRobustCommand{\DIFdelendFL}{\DIFOaddendFL \let\includegraphics\DIFOincludegraphics} %DIF PREAMBLE
%DIF AMSMATHULEM PREAMBLE %DIF PREAMBLE
\makeatletter %DIF PREAMBLE
\let\sout@orig\sout %DIF PREAMBLE
\renewcommand{\sout}[1]{\ifmmode\text{\sout@orig{\ensuremath{#1}}}\else\sout@orig{#1}\fi} %DIF PREAMBLE
\makeatother %DIF PREAMBLE
%DIF COLORLISTINGS PREAMBLE %DIF PREAMBLE
\RequirePackage{listings} %DIF PREAMBLE
\RequirePackage{color} %DIF PREAMBLE
\lstdefinelanguage{DIFcode}{ %DIF PREAMBLE
%DIF DIFCODE_UNDERLINE %DIF PREAMBLE
  moredelim=[il][\color{red}\sout]{\%DIF\ <\ }, %DIF PREAMBLE
  moredelim=[il][\color{blue}\uwave]{\%DIF\ >\ } %DIF PREAMBLE
} %DIF PREAMBLE
\lstdefinestyle{DIFverbatimstyle}{ %DIF PREAMBLE
	language=DIFcode, %DIF PREAMBLE
	basicstyle=\ttfamily, %DIF PREAMBLE
	columns=fullflexible, %DIF PREAMBLE
	keepspaces=true %DIF PREAMBLE
} %DIF PREAMBLE
\lstnewenvironment{DIFverbatim}[1][]{\lstset{style=DIFverbatimstyle}}{} %DIF PREAMBLE
\lstnewenvironment{DIFverbatim*}[1][]{\lstset{style=DIFverbatimstyle,showspaces=true}}{} %DIF PREAMBLE
\lstset{extendedchars=\true,inputencoding=utf8}

%DIF END PREAMBLE EXTENSION ADDED BY LATEXDIFF

\begin{document}
\vspace*{0.2in}

% Title must be 250 characters or less.
\begin{flushleft}
{\Large
\textbf{Evaluation of validation strategies and transferability of EEG attention decoding: a comparative analysis on open datasets} % Please use "sentence case" for title and headings.
}
\newline
% Insert author names, affiliations and corresponding author email (do not include titles, positions, or degrees).
\\
Sanzhar Sovet\textsuperscript{1},
Beibit Abdikenov\textsuperscript{1*},
Manzura Zholdasova\textsuperscript{2,3},
Altyngul Kamzanova\textsuperscript{2},
Medet Mukushev\textsuperscript{1},
Almira Kustubayeva\textsuperscript{2,3}
\\
\bigskip
\textbf{1} Science and Innovation Center ``Artificial Intelligence'', Astana IT University, Astana, Kazakhstan
\\
\textbf{2} Department of Biophysics, Biomedicine, and Neuroscience, Al-Farabi Kazakh National University, Almaty, Kazakhstan
\\
\textbf{3} Brain Institute, Al-Farabi Kazakh National University, Almaty, Kazakhstan
\\
\bigskip

% Use the asterisk to denote corresponding authorship and provide email address in note below.
* beibit.abdikenov@astanait.edu.kz

\end{flushleft}
% Please keep the abstract below 300 words

\section*{Abstract}
Despite growing interest, the development of robust EEG classification models remains a challenge. EEG signals are inherently noisy, non-stationary, and exhibit high inter-subject and inter-session variability. Furthermore, EEG-based models often fail to generalize across different recording sessions, subjects, or task contexts, which limits their real-world applicability. One of the key contributors to this limitation is the lack of rigorous validation procedures and over-reliance on subject-dependent or intra-session evaluations that tend to overestimate model performance. In this work, two publicly available datasets recorded with technically compatible Emotiv headsets but distinct cognitive paradigms were aligned through a unified binary attention labeling scheme, a shared subset of 12 electrodes, and a standardized preprocessing pipeline including bandpass filtering, independent component analysis, short-time Fourier transform, spectral band decomposition, and overlapping window segmentation. Model performance was assessed under within-dataset, cross-subject, cross-session, cross-task, and cross-dataset validation settings. Within-dataset evaluation yielded the strongest results, with accuracy and F1-score reaching 0.86 for Dataset A and 0.81 and 0.88, respectively, for Dataset B. Performance declined under more stringent transfer settings: cross-subject and cross-session validation showed moderate robustness, cross-task transfer was only partially successful when tasks shared cognitive demands, and cross-dataset transfer remained the most difficult scenario, with one transfer direction approaching chance level. These findings indicate that current EEG classifiers capture attention-related features that are strongly shaped by experimental context rather than invariant neural signatures of attention. The study highlights the importance of rigorous validation protocols and underscores the need for future methods that improve cross-paradigm generalization in EEG attention decoding.

\clearpage
\newgeometry{top=0.85in,left=1in,right=1in,footskip=0.75in}
\linenumbers

% Use "Eq" instead of "Equation" for equation citations.
\section*{Introduction}

Electroencephalography (EEG) has become a key modality in the development of brain--computer interface (BCI) systems and cognitive monitoring tools, due to its non-invasive nature, high temporal resolution, and ability to capture brain dynamics associated with cognitive states such as attention, workload, and emotional arousal. In particular, EEG-based attention recognition has gained increasing importance in applications ranging from neurofeedback and education to safety-critical systems such as drowsiness detection in transportation and adaptive human--machine interfaces.

Despite growing interest, the development of robust EEG classification models remains challenging. EEG signals are inherently noisy, non-stationary, and exhibit high inter-subject and inter-session variability. Furthermore, EEG-based models often fail to generalize across different recording sessions, subjects, or task contexts, which limits their real-world applicability. One of the key contributors to this limitation is the lack of rigorous validation procedures and over-reliance on subject-dependent or intra-session evaluations that tend to overestimate model performance.

In response to these issues, recent literature has emphasized the importance of evaluating EEG models under more realistic scenarios, including cross-subject, cross-session, and cross-task validation schemes. However, less attention has been paid to the problem of cross-paradigm transferability of neural markers, that is, whether EEG features learned in one task paradigm (e.g., passive observation with predefined attentional conditions) can generalize to another (e.g., active cognitive testing tasks). This question is particularly relevant for the development of generalizable attention-detection models.

\DIFaddbegin \DIFadd{\rev{A common route to these problems is transfer learning, whose principle is to reuse knowledge acquired on completed tasks in order to support related new ones~\cite{zhang2020transfer}. Its success depends not only on the modeling and learning strategy but also on the degree of domain and label compatibility between source and target data~\cite{wu2022transfer}. For cognitive constructs such as attention this compatibility is difficult to guarantee, because attention lacks a single widely accepted definition~\cite{hommel2019noone,sun2025attention}. Datasets designed to assess attention may therefore probe different attentional types, such as selective, sustained, alternating, divided, or focused attention, while assigning them the same nominal labels. The alignment between source and target datasets may vary accordingly, producing label mismatch and inconsistent decision boundaries under domain shift.}
}

\DIFaddend Another major challenge is that many comparative EEG studies suffer from technical mismatches between datasets, such as differences in hardware, number of channels, electrode placement, or sampling rates, which introduce distributional shifts and confound performance assessments. For example, prior domain adaptation studies have shown poor transfer performance due to such technical discrepancies rather than inherent limitations of neural representation.

To address these gaps, this study proposes a comparative analysis of EEG-based attention classification across two open-access datasets recorded with the same EEG device (Emotiv headset), but designed under distinct cognitive paradigms. One dataset uses explicit attention labels in a continuous, passive-control setting, while the other captures attention-related processes through well-established cognitive tasks, such as the Stroop test and mental arithmetic.

By aligning these datasets conceptually and technically, we aim to examine the transferability of neural attention markers across paradigms, and systematically evaluate how different validation strategies impact model generalization.

\DIFaddbegin \DIFadd{\rev{Prior work has emphasized that the efficacy of cross-dataset knowledge transfer and adaptation remains an open question in EEG decoding~\cite{lan2019domain}, and recent surveys of cross-domain EEG modeling reach the same conclusion~\cite{li2025foundation,verma2023attentionet}. The present analysis addresses this gap by quantifying zero-shot generalization in EEG-based attention decoding across cross-subject, cross-session, cross-task, and cross-dataset (cross-paradigm) settings, thereby establishing a baseline against which transfer-learning and domain-adaptation methods can be compared.}
}

\DIFaddend \section*{Related work}

The development of EEG-based classification models that generalize across subjects, sessions, and tasks is an active area of research in brain--computer interface (BCI) and cognitive monitoring fields. A growing body of work has explored how different validation strategies affect the robustness and transferability of EEG models. This section reviews relevant studies on cross-validation, domain adaptation, and cross-task EEG analysis, with a specific focus on attention and cognitive state recognition.

One of the most prominent lines of research involves domain adaptation techniques to reduce the discrepancy between training and testing distributions. For instance,~\DIFdelbegin \DIFdel{\mbox{%DIFAUXCMD
\cite{li2020domain} }\hskip0pt%DIFAUXCMD
}\DIFdelend \DIFaddbegin \DIFadd{\mbox{%DIFAUXCMD
\cite{lan2019domain} }\hskip0pt%DIFAUXCMD
}\DIFaddend conducted a comprehensive evaluation of feature-based domain adaptation methods, including MIDA, TCA, GFK, and others, on two widely used emotion-related EEG datasets: DEAP and SEED. Their experiments included leave-one-subject-out and cross-dataset validation, where models were trained on one dataset and tested on the other. The results revealed limited improvements, which the authors attributed to substantial differences in recording hardware and data distributions between the datasets. This finding underscores the importance of technical consistency in cross-dataset evaluations, a factor addressed in the present study, where both datasets were collected using Emotiv EEG headsets with comparable specifications.

\cite{chen2025hybrid} proposed a hybrid feature extraction approach for classifying mental attention states under cross-session and cross-subject scenarios. Their method combines spectral features obtained via short-time Fourier transform (STFT) with measures of functional and effective brain connectivity and employs a two-stage feature selection strategy to enhance model generalization. Experiments on two publicly available datasets demonstrated high classification accuracy (up to 94.01\% in cross-subject evaluation), confirming the potential of integrating spectral and network-based features to improve the robustness of attention decoding. \DIFaddbegin \DIFadd{\rev{Those datasets were \textit{EEG data for Mental Attention State Detection}~\cite{aci2019attention} and MEMA; the former is the dataset referred to as Dataset~A throughout the present study, so their cross-subject figures and ours are obtained on the same recordings.} }\DIFaddend However, this work focuses on cross-session and cross-subject scenarios and does not address the transferability of models across different tasks, datasets, or paradigms. The present study is aimed at filling this gap by examining the extent to which attention-related neural markers retain their discriminative power when transferred across cognitive paradigms.

Several studies have investigated the cross-task generalization of EEG features. In a notable example,~\DIFdelbegin \DIFdel{\mbox{%DIFAUXCMD
\cite{zhang2021cross} }\hskip0pt%DIFAUXCMD
}\DIFdelend \DIFaddbegin \DIFadd{\mbox{%DIFAUXCMD
\cite{ke2021crosstask} }\hskip0pt%DIFAUXCMD
}\DIFaddend analyzed the cross-task consistency of EEG-based mental workload indicators using a custom dataset involving two distinct paradigms: the N-back task and the Multi-Attribute Task Battery (MATB). The authors compared the performance of classifiers trained and tested within the same task versus across tasks. Their findings suggested that while task-specific models performed better, certain EEG features, especially in the power spectral domain, exhibited moderate stability across task contexts.

Task generalization has also been studied in attention recognition. For example, \DIFdelbegin \DIFdel{Nguyen }\DIFdelend \DIFaddbegin \DIFadd{He }\DIFaddend et al.~\DIFdelbegin \DIFdel{\mbox{%DIFAUXCMD
\cite{nguyen2022kalman} }\hskip0pt%DIFAUXCMD
}\DIFdelend \DIFaddbegin \DIFadd{\mbox{%DIFAUXCMD
\cite{he2025kalman} }\hskip0pt%DIFAUXCMD
}\DIFaddend proposed a Kalman filtering-based model and evaluated its robustness across multiple datasets\DIFaddbegin \DIFadd{\rev{, among them the same \textit{EEG data for Mental Attention State Detection} recordings used here as Dataset~A}}\DIFaddend . Their results demonstrated the feasibility of cross-dataset attention classification, although model performance degraded when shifting between tasks with differing structures.

\rev{Beyond feature-based alignment, a growing literature tackles domain shift in EEG with adversarial and deep domain-adaptation frameworks (e.g., domain-adversarial training with gradient reversal and adversarial discriminative adaptation), as well as personalization and few-shot calibration strategies that fine-tune a shared model to each new subject. Representation-learning approaches based on transformer/patching architectures and self-supervised or masked pretraining have likewise reported improved within- and cross-subject performance on cognitive EEG datasets. These methods are complementary to the present work: rather than proposing a new adaptation algorithm, we quantify how far conventional spectral features and a standard classifier transfer under increasingly stringent validation, thereby establishing a baseline against which such specialized methods can be compared.}

While these studies highlight the potential for EEG-based cross-task transfer, they often suffer from two key limitations. First, many rely on custom or proprietary datasets, which reduces reproducibility and hinders broader benchmarking. Second, differences in data acquisition hardware often confound comparisons and make it difficult to isolate the true effect of task variability.

In contrast, our study contributes to this line of work by focusing on attention-related EEG datasets that are both publicly available and technically compatible. We specifically investigate whether attention-related neural patterns, learned from data collected in one paradigm (passive observation with explicit labeling), can generalize to a different paradigm (active cognitive testing), without the confounding effects of sensor mismatch. \rev{To the best of our knowledge, comparative evaluations of cross-paradigm transferability that use open EEG datasets recorded with the same device family remain scarce, and the present analysis is intended as a step toward filling this gap rather than an exhaustive benchmark.}

\section*{Open-access datasets}

\subsection*{EEG data for mental attention state detection (Dataset A)}

The dataset described in~\cite{aci2019attention} comprises \textbf{34 EEG recordings} collected from \textbf{five participants} over the course of \textbf{seven days}, with the fifth participant completing only \textbf{six sessions}. During each session, participants operated a virtual train using the \textit{Microsoft Train Simulator}. Each recording lasted approximately \textbf{35 to 55 minutes}. All experiments in this study were performed with healthy volunteer participants chosen among students.

To ensure consistency, all sessions were conducted on the same predefined route, designed to require minimal interaction. This setup constitutes a \textbf{low-intensity control task}, minimizing variability due to task demands and aligning with \textit{naturalistic stimulus-based protocols} as described in~\cite{chen2025hybrid}. Each session was divided into \textbf{three sequential phases}, each designed to reflect a different target mental state: \textbf{\textit{focused}}, \textbf{\textit{unfocused}}, and \textbf{\textit{drowsy}}. The participants simulated these mental states based on standardized instructions provided by the experiment supervisor:

\begin{itemize}
  \item \textbf{Focused (0--10 min):} Participants actively monitored and mentally engaged with the simulator, maintaining sustained attention to the virtual controls and screen events.
  \item \textbf{Unfocused (10--20 min):} Participants disengaged from the task, ceasing interaction and intentionally avoiding attentional focus, while being instructed to keep their eyes open and remain awake.
  \item \textbf{Drowsy (20--30 min):} Participants were allowed to relax, close their eyes, and enter a drowsy or lightly dozing state, simulating reduced cognitive arousal.
\end{itemize}

EEG signals were recorded using a \textbf{modified Emotiv headset}, equipped with wet electrodes and supporting wireless Bluetooth transmission. The device recorded data at a \textbf{sampling rate of 128~Hz} and with a \textbf{voltage resolution of 0.51~$\mu$V}. The dataset includes \textbf{14 EEG channels}, labeled according to the 10--20 international system: AF3, F7, F3, FC5, T7, P7, O1, O2, P8, T8, FC6, F4, F8, and AF4.

\subsection*{SAM-40: dataset of 40-subject EEG recordings (Dataset B)}

The \textbf{SAM-40} dataset, described in~\cite{ghosh2022sam40}, contains EEG recordings from 40 healthy participants (14 females and 26 males) with a mean age of 21.5 years. The purpose of data collection was to investigate short-term stress responses induced by the performance of various cognitive tasks, which correspond to a \textit{continuous cognitive engagement} paradigm as described in~\cite{chen2025hybrid}. The experimental protocol included four conditions:

\begin{itemize}
  \item \textbf{Resting state} --- EEG recording during relaxation while listening to music.
  \item \textbf{Stroop Color-Word Test} --- participants were required to identify the color of the ink in which a word was printed, while ignoring the word's semantic meaning. The task included two types of stimuli: \textit{congruent} (the color name matched the ink color) and \textit{incongruent} (the color name and ink color differed). Each execution of the task included 11 stimuli.
  \item \textbf{Arithmetic problem solving} --- participants mentally solved a mathematical problem and determined whether the presented answer was correct, responding with a ``thumbs up'' or ``thumbs down'' gesture. Each execution of the task included 6 arithmetic stimuli.
  \item \textbf{Mirror Image Recognition} --- participants were shown images and asked to determine whether they were mirror-symmetric or asymmetric, responding with the same gestures. Each execution of the task included 8 images.
\end{itemize}

Each task lasted 25 seconds, with 5-second rest intervals and a 10-second instruction period preceding the next task. After the completion of each trial, participants rated the level of stress experienced for each task on a 10-point scale. For each participant, three trials were recorded using different sets of stimuli.

EEG signals were recorded using an Emotiv Epoc Flex headset with 32 wet electrodes, labeled according to the international 10--20 system. Data were acquired at a sampling rate of 128~Hz and a voltage resolution of 0.51~$\mu$V. The 32 channels included in the dataset are: Cz, Fz, Fp1, F7, F3, FC1, C3, FC5, FT9, T7, CP5, CP1, P3, P7, PO9, O1, Pz, Oz, O2, PO10, P8, P4, CP2, CP6, T8, FT10, FC6, C4, FC2, F4, F8, and Fp2.

\section*{Materials and methods}

The methodological design of this study was developed to evaluate the generalizability and transferability of EEG-based attention decoding across different recording conditions, participants, and task paradigms. To this end, we adopted two publicly available datasets that differ in experimental protocols but share compatible technical specifications, which allows for a controlled comparative analysis. The datasets were harmonized through a unified labeling scheme that distinguishes between low-attention and high-attention states, enabling consistent evaluation across scenarios.

\subsection*{Dataset alignment and integration rationale}

Effective evaluation of cross-task and cross-dataset generalization relies on the availability of datasets that are compatible both technically and conceptually. In this study, we integrate two publicly available EEG datasets that differ in their original purpose and labeling schemes, yet share critical properties enabling a unified analysis of attentional states.

A major strength of this dataset combination lies in their acquisition setup. Both datasets were recorded using consumer-grade Emotiv EEG headsets with identical electrode configurations, sampling rate (128~Hz), and voltage resolution. This uniformity eliminates a key source of variability, namely hardware-induced distribution shifts, that frequently hinder cross-dataset transferability as shown by~\DIFdelbegin \DIFdel{\mbox{%DIFAUXCMD
\cite{li2020domain}}\hskip0pt%DIFAUXCMD
}\DIFdelend \DIFaddbegin \DIFadd{\mbox{%DIFAUXCMD
\cite{lan2019domain}}\hskip0pt%DIFAUXCMD
}\DIFaddend . Consequently, performance differences observed in our experiments can be attributed primarily to task and participant factors, rather than to mismatched recording systems.

Although the first dataset contains explicit attention state labels (\textit{Focused}, \textit{Unfocused}, \textit{Drowsy}) collected during a low-interaction simulation task, and the second was originally designed for stress assessment, both capture cognitive conditions strongly tied to attentional processes. The second dataset includes tasks widely recognized in cognitive neuroscience for eliciting high levels of attentional engagement:

\begin{itemize}
  \item \textbf{Stroop Color Test} --- probes both \textit{selective} and \textit{sustained attention}, requiring participants to suppress automatic reading responses~\DIFdelbegin \DIFdel{\mbox{%DIFAUXCMD
\cite{macleod1991stroop,scarpina2017stroop,esterman2013neural}}\hskip0pt%DIFAUXCMD
}\DIFdelend \DIFaddbegin \DIFadd{\mbox{%DIFAUXCMD
\cite{macleod1991stroop,scarpina2017stroop,esterman2013zone}}\hskip0pt%DIFAUXCMD
}\DIFaddend .
  \item \textbf{Mental arithmetic task} --- demands \textit{sustained internal attention} and working memory for continuous manipulation of numerical information~\DIFdelbegin \DIFdel{\mbox{%DIFAUXCMD
\cite{grabner2007math,zhou2010eeg}}\hskip0pt%DIFAUXCMD
}\DIFdelend \DIFaddbegin \DIFadd{\mbox{%DIFAUXCMD
\cite{grabner2007math,kitaura2017functional}}\hskip0pt%DIFAUXCMD
}\DIFaddend .
  \item \textbf{Mirror image recognition task} --- engages visuospatial processing and requires both \textit{selective} and \textit{sustained attention} to discriminate between symmetric and asymmetric patterns~\DIFdelbegin \DIFdel{\mbox{%DIFAUXCMD
\cite{shepard1971mental,hyun2007visual,voyer2011performance}}\hskip0pt%DIFAUXCMD
}\DIFdelend \DIFaddbegin \DIFadd{\mbox{%DIFAUXCMD
\cite{shepard1971mental,hyun2007visual,tomasino2016mental}}\hskip0pt%DIFAUXCMD
}\DIFaddend .
\end{itemize}

These two datasets correspond to the \textit{naturalistic stimulus-based protocols} and \textit{continuous cognitive engagement} paradigms described in~\cite{chen2025hybrid}. In this unified labeling scheme, conditions such as \textit{Rest} or \textit{Unfocused} are treated as low-attention states, whereas \textit{Focused}, \textit{Stroop}, \textit{Arithmetic}, and \textit{Mirror recognition} are treated as high-attention states. \rev{The \textit{Drowsy} condition of Dataset~A denotes a state of reduced cognitive arousal and disengagement from the task; it was therefore mapped to the low-attention class together with \textit{Unfocused} and \textit{Rest}, rather than being discarded.} This mapping is summarized in Table~\ref{tab:label_mapping}.

\begin{table}[!ht]
\centering
\caption{\textbf{Unified binary attention label mapping across datasets.}}
\label{tab:label_mapping}
\begin{tabular}{|l|l|}
\hline
\textbf{Attention state} & \textbf{Conditions} \\ \hline
High attention & Focused, Stroop, Arithmetic, Mirror \\ \hline
Low attention  & Unfocused, Rest, \DIFdelbeginFL %DIFDELCMD < \cellcolor{yellow}%%%
\DIFdelendFL Drowsy \\ \hline
\end{tabular}
\end{table}

\subsection*{Experimental design}

The experimental design was structured to capture different dimensions of model generalization, ranging from subject- and session-level variability to cross-task and cross-paradigm transferability. Instead of relying on a single evaluation scheme, we combined five complementary validation strategies, each addressing a distinct source of variability.

\textbf{Baseline (within-dataset evaluation).}
Within-dataset performance was estimated using stratified K-fold cross-validation applied separately to each dataset. This setting provides an upper bound of model performance when training and testing distributions are identical, serving as a reference for subsequent generalization analyses. \rev{Unless otherwise stated, K was set to 5. Because the continuous recordings were segmented into densely overlapping windows (see Data preparation), windows derived from the same recording are temporally autocorrelated; we discuss the resulting risk of optimistic within-dataset estimates in the Limitations subsection.}

\textbf{Cross-subject validation.}
Generalization across individuals was examined using a leave-one-subject-out (LOSO) protocol. In each iteration, data from one participant were held out for testing, while the model was trained on all remaining participants. This evaluation reflects the model's ability to adapt to previously unseen subjects.

\textbf{Cross-session validation.}
To test robustness against intra-individual variability (e.g., electrode placement shifts, fatigue, or day-to-day differences), a leave-one-session-out scheme was adopted for the multi-session dataset. Here, the model was trained on all but one recording session of a participant and tested on the held-out session.

\textbf{Cross-task validation.}
Cross-task generalization was assessed within Dataset~B, which includes multiple cognitive tasks (Stroop, Arithmetic, and Mirror recognition). The model was trained exclusively on one task and evaluated on another (e.g., Stroop $\rightarrow$ Arithmetic). This protocol probes whether neural markers of attention learned in one cognitive paradigm retain discriminative power in a different paradigm. Similar cross-task transfer analyses have been conducted in prior studies on EEG workload and attention~\DIFdelbegin \DIFdel{\mbox{%DIFAUXCMD
\cite{zhang2021cross}}\hskip0pt%DIFAUXCMD
}\DIFdelend \DIFaddbegin \DIFadd{\mbox{%DIFAUXCMD
\cite{ke2021crosstask}}\hskip0pt%DIFAUXCMD
}\DIFaddend . \rev{In the present cross-task experiments the source and target tasks were drawn from the same pool of Dataset~B participants; consequently, task transfer and subject identity are not fully disentangled, a confound we revisit in the Limitations subsection.}

\textbf{Cross-dataset validation.}
Finally, cross-dataset transferability was investigated by training models on one entire dataset and testing them on the other (Dataset~A $\rightarrow$ Dataset~B, and vice versa). Additionally, subset experiments were performed (e.g., training on Dataset~A and testing only on Stroop trials of Dataset~B). This protocol can also be interpreted as a form of \textit{cross-paradigm} evaluation in the taxonomy proposed by~\cite{chen2025hybrid}. \DIFaddbegin \DIFadd{\rev{In that taxonomy the train-simulator protocol of Dataset~A falls under \textit{naturalistic stimulus-based} paradigms, whereas the neurocognitive tests of Dataset~B fall under \textit{continuous cognitive engagement}; the two datasets therefore differ at the paradigm level and not only at the dataset level.} }\DIFaddend The approach follows established cross-dataset validation protocols used in EEG research~\DIFdelbegin \DIFdel{\mbox{%DIFAUXCMD
\cite{li2020domain,nguyen2022kalman}}\hskip0pt%DIFAUXCMD
}\DIFdelend \DIFaddbegin \DIFadd{\mbox{%DIFAUXCMD
\cite{lan2019domain,he2025kalman}}\hskip0pt%DIFAUXCMD
}\DIFaddend .

For performance evaluation, we report four widely used classification metrics: \textit{Accuracy}, \textit{Precision}, \textit{Recall}, and \textit{F1-score}. To ensure fair assessment in the presence of class imbalance, particularly in Dataset~B, \textit{Precision}, \textit{Recall}, and \textit{F1-score} were computed using a weighted averaging scheme.

\subsection*{Data preparation}

As the two datasets were recorded with different electrode montages (14 vs.\ 32 channels), only the common subset of 12 electrodes was retained for analysis: F3, F4, F7, F8, FC5, FC6, O1, O2, P7, P8, T7, and T8. Restricting the analysis to this overlapping set ensured consistency in the spatial representation of EEG signals and eliminated potential biases due to non-matching channel configurations.

Following channel selection, all EEG recordings were preprocessed using a standardized pipeline. Signals were bandpass-filtered in the range of 0.5--45~Hz to suppress line noise. Independent Component Analysis (ICA) was subsequently applied to identify and remove components associated with ocular artifacts. \rev{Components were flagged for removal based on their scalp topographies and time courses characteristic of eye-blink and lateral-eye-movement activity (frontal, low-frequency, high-amplitude patterns); typically two such components were excluded per recording.} \rev{Given that only 12 channels were retained, this low-density configuration limits the stability of ICA decomposition; we note this constraint in the Limitations subsection.}

After preprocessing, the EEG signals were transformed into the time--frequency domain using the Short-Time Fourier Transform (STFT). Instead of analyzing the entire broadband spectrum, the resulting spectrograms were decomposed into five canonical frequency bands most relevant for cognitive load and attentional processes: \textit{delta} (0.5--4~Hz), \textit{theta} (4--8~Hz), \textit{alpha} (8--13~Hz), \textit{beta} (13--30~Hz), and \textit{gamma} (30--45~Hz).

The continuous EEG recordings were segmented into overlapping windows of 1~second duration with a step size of 0.125~seconds. This windowing strategy approximates the requirements of real-time applications while maintaining sufficient temporal and spectral resolution for robust analysis.

\subsection*{\revs{Classification model and training configuration}}

\rev{Classification was performed with a gradient-boosted decision-tree model (XGBoost), which is well suited to the resulting tabular, high-dimensional band-power feature representation and has been reported to perform robustly in EEG decoding tasks. Per-band, per-channel STFT power features were standardized prior to training. The same model family and feature representation were used across all validation scenarios so that differences in performance reflect the validation protocol rather than changes in the learner. The complete preprocessing and evaluation pipeline, including the exact model hyperparameters and data splits, is publicly available at \url{https://github.com/sovetskiysn/eeg-validation-strategies}.}

% Place figure captions after the first paragraph in which they are cited.
%DIF >  PLOS ONE requires each figure to be uploaded as an individual file (Fig1.tif,
%DIF >  Fig2.tif) separate from the manuscript; only the caption belongs here.
%DIF >  Do not add \includegraphics to this file.
The overall decay of performance as validation becomes more stringent is illustrated in Fig~\ref{fig:performance_decay}, and the full set of pairwise transfer outcomes within Dataset~B is shown in Fig~\ref{fig:cross-task}.

\DIFdelbegin %DIFDELCMD < \begin{figure}[!ht]
%DIFDELCMD < \centering
%DIFDELCMD < \includegraphics[width=\linewidth]{figures/fig1.jpg}
%DIFDELCMD < %%%
\DIFdelendFL \DIFaddbeginFL \begin{figure}[!h]
\DIFaddendFL \caption{\textbf{Performance decay across validation strategies.}
Classification performance decreases as the validation protocol moves from within-dataset evaluation toward cross-subject, cross-session, cross-task, and cross-dataset transfer.}
\label{fig:performance_decay}
\end{figure}

\DIFdelbegin %DIFDELCMD < \begin{figure}[!ht]
%DIFDELCMD < \centering
%DIFDELCMD < \includegraphics[width=0.72\linewidth]{figures/fig2.jpg}
%DIFDELCMD < %%%
\DIFdelendFL \DIFaddbeginFL \begin{figure}[!h]
\DIFaddendFL \caption{\textbf{Cross-task generalization matrix.}
Pairwise transfer performance between the cognitive tasks of Dataset~B (Stroop, Arithmetic, and Mirror recognition).}
\label{fig:cross-task}
\end{figure}

\section*{Results and discussion}

All experimental results are summarized in Table~\ref{tab:results}. The dataset introduced in~\cite{aci2019attention} is denoted as \textit{Full A}, while the dataset described in~\cite{ghosh2022sam40} is denoted as \textit{Full B}. This unified presentation facilitates direct comparison of validation schemes and transfer scenarios across both datasets.

The evaluation revealed a consistent degradation of performance as validation protocols became more stringent. Within-dataset validation (stratified K-fold) yielded the highest scores, with Dataset A reaching 0.86 in all metrics, and Dataset B performing comparably with 0.81 accuracy and 0.88 F1-score. Notably, when Dataset B tasks were considered separately, performance was lower for Stroop (0.78 F1), Arithmetic (0.72 F1), and Mirror (0.76 F1). However, when combined into a unified dataset, the model achieved superior overall performance, particularly in recall (0.97). \rev{This high recall, paired with a lower precision and F1-score, is consistent with a class imbalance in the harmonized Dataset~B and a decision boundary biased toward the majority (high-attention) class; reporting of class distributions and uncertainty is addressed as a limitation below.} This suggests that task aggregation amplifies the salience of attentional markers by capturing complementary neural mechanisms across paradigms, resulting in stronger discriminability than task-specific decoders.

% Place tables after the first paragraph in which they are cited.
\begin{table}[!ht]
\centering
\caption{\textbf{Performance metrics across all experimental scenarios.}}
\label{tab:results}
\begin{adjustwidth}{-0.3in}{-0.3in}
\centering
\begin{tabular}{|l|l|c|c|c|c|}
\hline
\textbf{Scenario} & \textbf{Validation method} & \textbf{Accuracy} & \textbf{Precision} & \textbf{Recall} & \textbf{F1-score} \\ \hline
Baseline (within) & Stratified K-fold (Full A)             & 0.86 & 0.86 & 0.86 & 0.86 \\ \hline
Baseline (within) & Stratified K-fold (Full B)             & 0.81 & 0.81 & 0.97 & 0.88 \\ \hline
Baseline (within) & Stratified K-fold (B: Stroop only)     & 0.78 & 0.78 & 0.79 & 0.78 \\ \hline
Baseline (within) & Stratified K-fold (B: Arithmetic only) & 0.72 & 0.71 & 0.72 & 0.72 \\ \hline
Baseline (within) & Stratified K-fold (B: Mirror only)     & 0.75 & 0.75 & 0.77 & 0.76 \\ \hline
Cross-subject & Leave-one-subject-out (Full A)             & 0.63 & 0.61 & 0.68 & 0.64 \\ \hline
Cross-subject & Leave-one-subject-out (Full B)             & 0.73 & 0.77 & 0.90 & 0.83 \\ \hline
Cross-session & Leave-one-session-out (Full A)             & 0.64 & 0.63 & 0.71 & 0.66 \\ \hline
Cross-session & Leave-one-session-out (Full B)             & 0.76 & 0.78 & 0.94 & 0.85 \\ \hline
Cross-task (only B) & Stroop $\rightarrow$ Arithmetic      & 0.57 & 0.58 & 0.57 & 0.56 \\ \hline
Cross-task (only B) & Stroop $\rightarrow$ Mirror          & 0.70 & 0.80 & 0.70 & 0.67 \\ \hline
Cross-task (only B) & Arithmetic $\rightarrow$ Stroop      & 0.59 & 0.60 & 0.59 & 0.59 \\ \hline
Cross-task (only B) & Arithmetic $\rightarrow$ Mirror      & 0.54 & 0.54 & 0.54 & 0.53 \\ \hline
Cross-task (only B) & Mirror $\rightarrow$ Stroop          & 0.78 & 0.84 & 0.78 & 0.78 \\ \hline
Cross-task (only B) & Mirror $\rightarrow$ Arithmetic      & 0.52 & 0.57 & 0.52 & 0.43 \\ \hline
Cross-dataset & Full A $\rightarrow$ Full B                & 0.75 & 0.73 & 0.75 & 0.64 \\ \hline
Cross-dataset & Full B $\rightarrow$ Full A                & 0.49 & 0.49 & 0.49 & 0.48 \\ \hline
Cross-dataset & Full A $\rightarrow$ Stroop only           & 0.50 & 0.67 & 0.50 & 0.34 \\ \hline
Cross-dataset & Full A $\rightarrow$ Arithmetic only       & 0.50 & 0.64 & 0.50 & 0.33 \\ \hline
Cross-dataset & Full A $\rightarrow$ Mirror only           & 0.50 & 0.56 & 0.50 & 0.34 \\ \hline
Cross-dataset & Stroop only $\rightarrow$ Full A           & 0.50 & 0.50 & 0.50 & 0.50 \\ \hline
Cross-dataset & Arithmetic only $\rightarrow$ Full A       & 0.50 & 0.50 & 0.50 & 0.45 \\ \hline
Cross-dataset & Mirror only $\rightarrow$ Full A           & 0.48 & 0.47 & 0.48 & 0.43 \\ \hline
\end{tabular}
\begin{flushleft} Precision, Recall, and F1-score are reported using weighted averaging. Arrows ($\rightarrow$) denote the direction of transfer (train $\rightarrow$ test).
\end{flushleft}
\end{adjustwidth}
\end{table}

In cross-subject and cross-session evaluations, accuracy dropped to 0.63--0.64 for Dataset A and 0.73--0.76 for Dataset B, reflecting the well-documented challenges of inter-individual variability and session non-stationarity in EEG data. Despite this degradation, Dataset B maintained relatively high F1-scores (up to 0.85), suggesting that task-induced attentional engagement provides more transferable features across individuals. These findings align with previous studies emphasizing the fragility of within-subject models when exposed to distributional shifts and highlight the potential of task-driven paradigms for mitigating such variability~\DIFdelbegin \DIFdel{\mbox{%DIFAUXCMD
\cite{chen2025hybrid,abdulridha2023attention}}\hskip0pt%DIFAUXCMD
.
}\DIFdelend \DIFaddbegin \DIFadd{\mbox{%DIFAUXCMD
\cite{chen2025hybrid}}\hskip0pt%DIFAUXCMD
.
%DIF >  TODO(authors) -- BLOCKER: the submitted version also cited the key
%DIF >  `abdulridha2023attention` here, but no BibTeX entry for it exists anywhere in
%DIF >  the project. A literature search during the round-1 bibliography check did
%DIF >  not surface any EEG attention paper by an author named Abdulridha, so the
%DIF >  intended source remains unidentified. The key stays dropped so that the
%DIF >  manuscript compiles without an unresolved "[?]" citation. Restore it as
%DIF >  \cite{chen2025hybrid,abdulridha2023attention} only once the intended
%DIF >  reference is identified and added to refs.bib.
}\DIFaddend 

Cross-task transfers showed that generalization strongly depends on the cognitive similarity of tasks. Transfers between Stroop and Mirror achieved moderate success (up to 0.78 accuracy), suggesting that these paradigms engage partially overlapping attentional mechanisms. In contrast, transfers involving Arithmetic were consistently weaker (0.54--0.59 accuracy). This discrepancy likely arises from the substantial differences between the neurocognitive tasks: while Stroop and Mirror share elements of selective attention and visual discrimination, Arithmetic introduces a fundamentally distinct mode of processing. These results emphasize that the heterogeneity of task demands directly impacts the stability of attention-related neural markers and limits their transferability.

Cross-dataset transfers remained the most challenging scenario. Training on Dataset A and testing on Dataset B yielded moderate accuracy (0.75), but the relatively low F1-score (0.64) indicates a bias toward the majority class, suggesting limited extraction of task-invariant attention features. Conversely, transfers in the opposite direction collapsed to chance-level performance (0.49--0.50 across metrics), underscoring the strong dataset-specific nature of EEG representations. These findings further demonstrate that while cross-subject, cross-session, and even certain cross-task evaluations suggest partial preservation of attention-related neural markers, such markers do not necessarily generalize across datasets. The near-random performance observed in cross-dataset settings highlights the fragility of these features when confronted with paradigm-level variability.

\rev{The cross-dataset results reported here characterize the difficulty of the problem under a conventional pipeline; they do not represent an upper bound. Dedicated domain-adaptation methods, including correlation alignment (CORAL), transfer component analysis (TCA), domain-adversarial training with gradient reversal (DANN), adversarial discriminative domain adaptation (ADDA), and Riemannian alignment, as well as few-shot per-subject calibration, are explicitly designed to mitigate the distribution shift observed here and could plausibly raise cross-paradigm performance above chance. Because such adaptation and personalization baselines were not implemented in this study, the present cross-dataset numbers should be read as a baseline lower bound rather than evidence that the transfer problem is intractable.}

Taken together, these results illustrate that while EEG-based models can reliably capture attentional states within a single dataset, their performance does not extend robustly across subjects, sessions, tasks, or paradigms. The observed failures in transferability emphasize that current classifiers exploit features tied to specific experimental conditions rather than invariant neural signatures of attention. Addressing this limitation will require methodological advances capable of disentangling paradigm-specific effects from more generalizable neural markers.

\subsection*{\revs{Limitations and future work}}

\rev{Several limitations qualify the present findings and motivate future work. First, the within-dataset baselines were computed over densely overlapping 1-second windows (0.125-second step); windows from the same recording are highly autocorrelated, so stratified K-fold that splits such windows across folds can leak temporally adjacent samples and inflate within-dataset metrics. Group-aware splitting by trial, session, or subject, or temporal blocking, would yield more conservative and realistic estimates. Second, all reported numbers are point estimates: class distributions after label harmonization, fold-wise variability, confidence intervals, and statistical significance tests are not reported, which limits the strength of quantitative comparisons. Third, in the cross-task experiments the source and target tasks share the same Dataset~B participants, so task transfer is not fully separated from subject identity; subject-disjoint cross-task splits are needed to isolate the task effect. Fourth, the low-density 12-channel montage constrains the stability of ICA-based artifact removal. Finally, the study does not include domain-adaptation or personalization baselines or additional diagnostic visualizations (e.g., confusion matrices, per-subject variability, or low-dimensional embeddings of the feature space). Addressing these points, releasing standardized leakage-free splits, and benchmarking adaptation methods are natural next steps that would strengthen the evidential value of this analysis.}

\section*{Conclusion}

This work systematically evaluated EEG-based attention classification under multiple validation schemes, ranging from within-dataset baselines to the most challenging cross-dataset transfers. The results reveal a consistent pattern: while within-dataset models achieve high accuracy, their robustness declines once exposed to new subjects, sessions, tasks, or paradigms. Cross-task generalization is only partially successful when cognitive mechanisms overlap, as in Stroop and Mirror, but fails for tasks that rely on fundamentally different processes such as arithmetic. Notably, although individual tasks in Dataset B yielded lower scores in isolation, their combination produced the strongest overall performance, suggesting that aggregating cognitively diverse paradigms enhances the salience of attention-related neural markers. Most critically, cross-dataset/cross-paradigm transfer remains an unsolved problem, as paradigm-specific variability introduces distributional shifts that substantially limit the effectiveness of conventional classifiers.

\DIFdelbegin \section*{\DIFdel{Funding}}
%DIFAUXCMD
\DIFdelend %DIF >  PLOS ONE ending sections, in the required order:
%DIF >    Acknowledgments -> References -> Supporting information captions
%DIF >  (https://journals.plos.org/plosone/s/submission-guidelines)
%DIF > 
%DIF >  The "Funding" section that appeared in the submitted version has been removed.
%DIF >  PLOS states: "Do not include funding sources in the Acknowledgments or anywhere
%DIF >  else in the manuscript file. Funding information should only be entered in the
%DIF >  financial disclosure section" of Editorial Manager. The funder, grant number
%DIF >  BR27198099 and the Role of Funder statement are carried in
%DIF >  submission_metadata.md instead.

\DIFdelbegin \DIFdel{This study was funded by the Committee of Science of the Ministry of Science and Higher Education of the Republic of Kazakhstan (Grant No. BR27198099).
}\DIFdelend %DIF >  TODO(authors): add an Acknowledgments section here only if there are
%DIF >  non-author contributors to thank. Do not list funding here. Leave the section
%DIF >  out entirely if there is nothing to acknowledge.
%DIF >  \section*{Acknowledgments}
%DIF >  ...

\nolinenumbers

\bibliography{refs}
\DIFaddbegin 

%DIF >  ---------------------------------------------------------------------------
%DIF >  TODO(authors) -- BLOCKER: Supporting information captions.
%DIF > 
%DIF >  PLOS ONE requires that every Supporting Information file have a caption
%DIF >  listed in the manuscript AFTER the reference list, and that the file be
%DIF >  uploaded in Editorial Manager with file type "Supporting Information"
%DIF >  (editor letter, item 6). The supplementary figures of PONE-D-26-28965 were
%DIF >  uploaded as type "Figure" and their source files are not present in this
%DIF >  repository, so their titles cannot be reconstructed here without inventing
%DIF >  content.
%DIF > 
%DIF >  Uncomment the block below and complete one entry per supporting file, one
%DIF >  "S<n> Fig." entry per uploaded supplementary figure, in the order they are
%DIF >  first cited. Required format (https://journals.plos.org/plosone/s/supporting-information):
%DIF >      \paragraph*{S1 Fig.}
%DIF >      \label{S1_Fig}
%DIF >      {\bf Title of the supporting figure.} Optional legend text.
%DIF >  The uploaded filename must match the caption label, e.g. S1_Fig.tif for
%DIF >  "S1 Fig." Cite each item in the body text with \nameref{S1_Fig}.
%DIF > 
%DIF >  \section*{Supporting information}
%DIF > 
%DIF >  \paragraph*{S1 Fig.}
%DIF >  \label{S1_Fig}
%DIF >  {\bf TITLE TO BE PROVIDED.} Legend to be provided.
%DIF >  ---------------------------------------------------------------------------
\DIFaddend 

\end{document}